The IOAA2021 logo is formed by the acronym IOAA, where the first letter is
represented by the silhouette of the building of the National Astronomical
Observatory (OAN) of Colombia, the oldest observatory in America. This
observatory is located in Bogota, where it was founded in 1803. The capital
city of Colombia is bordered by two famous hills, Monserrate and its neighbor
Guadalupe, which are icons of Bogota's cityscape that decorate the logo's
background.

\ioaafig{2021_TH_Q8-f1.pdf}

\begin{center}
Aerial view of Bogota City. Numbers show locations for the quoted places: 1 is
for OAN; 2 is for Guadalupe; and 3 is for Monserrate.
\end{center}

\begin{center}
\begin{tabular}{|c|c|c|c|}
\hline
Point & Latitude & Longitude & \begin{tabular}{@{}c@{}}Elevation\\(m.a.s.l)\end{tabular} \\
\hline
1 & $4^\circ 35' 53''$ N & $74^\circ 04' 37''$ W & 2607 \\
\hline
2 & $4^\circ 35' 30''$ N & $74^\circ 03' 15''$ W & 3296 \\
\hline
3 & $4^\circ 36' 18''$ N & $74^\circ 03' 19''$ W & 3100 \\
\hline
\end{tabular}
\end{center}

\begin{description}
\item[(8.1)] Estimate the distance (in km), between points 2 (Guadalupe) and 3
  (Monserrate). \hfill[3.0\,pt]
\item[(8.2)] Estimate the angular separation (in degrees) between Guadalupe (2)
  and Monserrate (3) as observed from the National Astronomical Observatory of
  Colombia (1). \hfill[6.0\,pt]
\item[(8.3)] From the OAN, on September 21 at 8:00 p.m.\ the Moon was observed
  towards the eastern hills (between Monserrate and Guadalupe). The measured
  ecliptic coordinates (longitude and latitude) of the Moon are shown in the
  table. Determine the equatorial coordinates of the Moon at the time of
  observation. \hfill[6.0\,pt]
\end{description}

\ioaafig{2021_TH_Q8-f2.pdf}
\ioaafig{2021_TH_Q8-f3.pdf}

\begin{center}
\begin{tabular}{|l|}
\hline
\textbf{Local Time: 8:00 p.m.} \\
\hline
$Az: +90^\circ 42' 59''\;/\;Alt: +19^\circ 01' 42''$ \\
\hline
$\lambda: +12^\circ 20' 16''\;/\;\beta: -04^\circ 24' 14''$ \\
\hline
\end{tabular}
\end{center}

Note: Azimuth measured from North to East.
